\section{Related Work}
\label{sec:related-work}

ASZ relates to prior work along four dimensions: tunable compression mechanisms, candidate evaluation, search-space decomposition, and cross-file configuration reuse. This section reviews these areas, complementing the detailed discussion of SZ3 and QoZ in Section~\ref{sec:baseline-configurations}.

\subsection{Configurations Tuned by Existing Compressors}
\label{sec:related-work:dimensions}

Using the terminology defined in Section~\ref{sec:pipeline:terminology}, existing prediction-based compressors expose six main types of configuration choices.

\begin{enumerate}
\item \emph{Prediction Stencils.}\par\noindent\mbox{\hspace{1em}}SZauto chooses between first- and second-order Lorenzo predictors~\cite{zhao2020significantly}. SZ3 considers Lorenzo prediction as well as linear and cubic interpolation stencils. QoZ may further divide a prediction round into two passes, allowing second-pass targets to use values reconstructed in the first pass as additional references. The resulting variants include a six-reference one-dimensional cubic stencil and stencils that draw references from multiple dimensions. GPU-based cuSZ-Hi instead provides a parallel multidimensional interpolation scheme that progresses from one-dimensional to three-dimensional interpolation within each level~\cite{wu2025boosting}.

\item \emph{Prediction Axes.}\par\noindent\mbox{\hspace{1em}}SZ3, QoZ, and CliZ~\cite{jian2024cliz} tune prediction axes. QoZ can freeze less smooth axes, while CliZ permutes and fuses dimensions to direct more predictions along smoother dimensions. The multidimensional scheme in cuSZ-Hi is isotropic and removes axis-order selection.

\item \emph{Prediction Coefficients.}\par\noindent\mbox{\hspace{1em}}SZ3 uses one fixed coefficient set for linear interpolation and one for cubic interpolation. QoZ provides an additional coefficient set for cubic interpolation. For these one-dimensional stencils, the compressors select among predefined coefficient sets rather than directly tuning individual coefficients. For multidimensional prediction, QoZ derives axis-blending weights from sampled prediction errors and assigns larger weights to axes with smaller errors.

\item \emph{Quantization.}\par\noindent\mbox{\hspace{1em}}SZauto tunes the number of quantization codes. QoZ uses different quantization intervals across prediction rounds, while Jin et al. tune the intervals across data partitions.

\item \emph{Lossless Compression.}\par\noindent\mbox{\hspace{1em}}CliZ classifies bins into groups encoded by separate Huffman trees according to domain-specific patterns. Liu et al.~\cite{liu2021improving} compare lossless backends for SZ2.1 and recommend Zstd~\cite{ycollet2021zstd} for its balance of compression ratio and compression/decompression speed on the evaluated datasets. ASZ fixes this backend to Huffman coding followed by Zstd rather than including backend selection in its tuning space.

\item \emph{Pipeline Composition.}\par\noindent\mbox{\hspace{1em}}A compressor may choose which prediction or transform methods to use and how its components are combined. SZ2~\cite{liang2018error} and SZauto select between Lorenzo and regression predictors at the block level. Liang et al.~\cite{liang2019significantly} and FAZ~\cite{liu2023faz} extend this selection to transform methods.
\end{enumerate}

\subsection{Candidate Evaluation Methods}
\label{sec:related-work:evaluation}

To compare the configuration choices reviewed above, existing evaluation methods use three types of criteria.

\begin{enumerate}
\item \emph{Compression-Based Evaluation.} SZ3, SZauto, QoZ, CliZ, and FAZ evaluate candidate configurations on sampled data. FRaZ~\cite{underwood2020fraz} and OptZConfig~\cite{underwood2022optzconfig} repeatedly run a compressor while searching for configurations that meet user-specified compression-ratio or quality requirements.

\item \emph{Compression-Result Estimation.} These methods estimate compression ratio or reconstruction quality without executing the complete compression process. For example, Liang et al. compress the coefficients produced by their transform-based predictor using either the SZ pipeline or embedded bit-plane coding. They estimate the resulting bit rates using discrete entropy for SZ compression and coefficient bit patterns for bit-plane coding. Similarly, Jin et al. use the discrete entropy of quantized prediction errors to approximate Huffman coding size. To account for subsequent lossless compression, they further adjust this approximation based on the fraction quantized to zero. They also derive quality metrics such as PSNR from the estimated reconstruction-error distribution and use the results to guide predictor and error-bound selection. ASZ likewise evaluates the first-order entropy $H_1$ of QBs and PSNR from reconstruction errors. Sections~\ref{sec:configuration-space:coefficient-selection} and~\ref{sec:configuration-space:coefficient-evaluation} detail how its computation methods and purpose differ from those of Jin et al.

\item \emph{Alternative Evaluation Criteria.} These methods select configurations using other criteria without explicitly estimating the final compression ratio or reconstruction quality. Common criteria include mean squared and mean absolute prediction errors. SZ2, SZauto, cuSZ-Hi, and QoZ use prediction-error criteria to select predictors or interpolation configurations. QoZ also uses mean squared prediction error to determine multidimensional interpolation weights.
\end{enumerate}

\subsection{Search Space Decomposition}
\label{sec:related-work:search}

Searching all combinations of multiple tuning dimensions is expensive. Compressors therefore divide the joint search space into smaller subspaces and search them separately. SZ3 first searches combinations of linear or cubic interpolation stencils and axis orders, then compares the best interpolation configuration with Lorenzo prediction. QoZ divides the prediction rounds into groups of three and tunes one interpolation configuration for each group.

\subsection{Cross-File Configuration Reuse}
\label{sec:related-work:reuse}

Existing compressors reuse different types of tuning results across files. FRaZ first tests the error bound selected for the preceding time step and reuses it when it still meets the target compression ratio. OptZConfig similarly supports a supplied initial guess and stops without further search when that guess satisfies the requested metric. CliZ takes an offline--online approach, separating pipeline selection from compression and reusing selected pipelines across climate fields or snapshots within the same model.

